% =================================================
% TikZ-рисунки к листу «Длина отрезка. Измерение и построение»
% =================================================

\tikzset{
  geom line/.style={draw=Primary,line width=1.05pt,line cap=round},
  geom segment/.style={draw=Secondary,line width=1.45pt,line cap=round},
  geom auxiliary/.style={draw=GrayLine,line width=0.7pt},
  geom point/.style={circle,fill=Accent,draw=white,line width=0.35pt,inner sep=2.15pt},
  geom label/.style={font=\small\bfseries,text=GrayText},
  geom caption/.style={font=\small\sffamily,text=GrayText,align=center},
  geom object label/.style={font=\small\bfseries\sffamily,text=Secondary},
  geom card/.style={draw=GrayLine,fill=white,rounded corners=1.5mm,line width=0.65pt},
  geom ruler/.style={draw=GrayText,line width=0.7pt},
  geom tick major/.style={draw=GrayText,line width=0.75pt},
  geom tick minor/.style={draw=GrayText,line width=0.45pt},
  geom answer/.style={font=\small\bfseries\sffamily,text=Secondary}
}

% Универсальная линейка от 0 до 8 см с миллиметровыми делениями.
\newcommand{\drawruler}[1][8]{%
  \draw[geom ruler] (0,0) rectangle (#1,0.45);
  \foreach \x in {0,...,#1}{
    \draw[geom tick major] (\x,0.45) -- (\x,0.05);
    \node[font=\scriptsize\sffamily,text=GrayText,below=1pt] at (\x,0) {\x};
  }
  \foreach \x in {0,...,#1}{
    \foreach \k in {1,...,9}{
      \pgfmathsetmacro{\xx}{\x+0.1*\k}
      \ifdim \xx pt < #1pt
        \pgfmathtruncatemacro{\isfive}{mod(\k,5)}
        \ifnum\isfive=0
          \draw[geom tick minor] (\xx,0.45) -- (\xx,0.18);
        \else
          \draw[geom tick minor] (\xx,0.45) -- (\xx,0.28);
        \fi
      \fi
    }
  }
  \node[font=\scriptsize\sffamily,text=GrayText,anchor=west] at (#1+0.15,0.13) {см};
}

\newcommand{\segmentlengthfigure}{%
  \begin{center}
    \begin{tikzpicture}[x=1cm,y=1cm]
      \coordinate (A) at (-3.2,0);
      \coordinate (B) at (3.2,0);
      \draw[geom segment] (A) -- (B);
      \node[geom point] at (A) {};
      \node[geom point] at (B) {};
      \node[geom label,above=3pt] at (A) {$A$};
      \node[geom label,above=3pt] at (B) {$B$};
      \node[geom object label,below=8pt] at (0,0) {$AB=6\text{ см }4\text{ мм}$};
    \end{tikzpicture}
  \end{center}
}

\newcommand{\correctmeasurementfigure}{%
  \begin{center}
    \begin{tikzpicture}[x=1cm,y=1cm]
      \node[geom card,minimum width=15.6cm,minimum height=3.35cm] at (0,0.55) {};
      \begin{scope}[shift={(-4,0)}]
        \drawruler[8]
        \coordinate (A) at (0,1.25);
        \coordinate (B) at (5.7,1.25);
        \draw[geom segment] (A)--(B);
        \node[geom point] at (A) {};
        \node[geom point] at (B) {};
        \node[geom label,above=2pt] at (A) {$A$};
        \node[geom label,above=2pt] at (B) {$B$};
        \draw[geom auxiliary,densely dashed] (A)--(0,0.45);
        \draw[geom auxiliary,densely dashed] (B)--(5.7,0.45);
      \end{scope}
      \node[geom caption] at (0,-0.85) {Начало отрезка совмещено с делением $0$. Длина $AB=5\text{ см }7\text{ мм}$.};
    \end{tikzpicture}
  \end{center}
}

\newcommand{\shiftedmeasurementfigure}{%
  \begin{center}
    \begin{tikzpicture}[x=1cm,y=1cm]
      \node[geom card,minimum width=15.6cm,minimum height=3.55cm] at (0,0.55) {};
      \begin{scope}[shift={(-4,0)}]
        \drawruler[8]
        \coordinate (C) at (1.2,1.25);
        \coordinate (D) at (6.8,1.25);
        \draw[geom segment] (C)--(D);
        \node[geom point] at (C) {};
        \node[geom point] at (D) {};
        \node[geom label,above=2pt] at (C) {$C$};
        \node[geom label,above=2pt] at (D) {$D$};
        \draw[geom auxiliary,densely dashed] (C)--(1.2,0.45);
        \draw[geom auxiliary,densely dashed] (D)--(6.8,0.45);
      \end{scope}
      \node[geom caption,text width=13.4cm] at (0,-0.75)
        {Отрезок начинается не у нуля: из показания правого конца нужно вычесть показание левого.};
    \end{tikzpicture}
  \end{center}
}

\newcommand{\unitsladderfigure}{%
  \begin{center}
    \begin{tikzpicture}[x=1cm,y=1cm,node distance=0.42cm]
      \node[geom card,minimum width=2.5cm,minimum height=1.2cm] (km) {\sffamily\bfseries км};
      \node[geom card,minimum width=2.5cm,minimum height=1.2cm,right=of km] (m) {\sffamily\bfseries м};
      \node[geom card,minimum width=2.5cm,minimum height=1.2cm,right=of m] (dm) {\sffamily\bfseries дм};
      \node[geom card,minimum width=2.5cm,minimum height=1.2cm,right=of dm] (cm) {\sffamily\bfseries см};
      \node[geom card,minimum width=2.5cm,minimum height=1.2cm,right=of cm] (mm) {\sffamily\bfseries мм};
      \draw[-{Stealth[length=2.4mm]},draw=Accent,line width=0.9pt] (km)--node[above,font=\scriptsize\sffamily]{$\times1000$}(m);
      \draw[-{Stealth[length=2.4mm]},draw=Accent,line width=0.9pt] (m)--node[above,font=\scriptsize\sffamily]{$\times10$}(dm);
      \draw[-{Stealth[length=2.4mm]},draw=Accent,line width=0.9pt] (dm)--node[above,font=\scriptsize\sffamily]{$\times10$}(cm);
      \draw[-{Stealth[length=2.4mm]},draw=Accent,line width=0.9pt] (cm)--node[above,font=\scriptsize\sffamily]{$\times10$}(mm);
      \node[geom caption,text width=13.6cm,below=7pt of dm]
        {При переходе к меньшей единице число увеличивается. В обратную сторону выполняют деление.};
    \end{tikzpicture}
  \end{center}
}

\newcommand{\constructionalgorithmfigure}{%
  \begin{center}
    \begin{tikzpicture}[x=1cm,y=1cm]
      \foreach \x/\title in {-5.2/{1. Отметьте $A$},0/{2. Отложите длину},5.2/{3. Отметьте $B$}}{
        \node[geom card,minimum width=4.7cm,minimum height=2.45cm] at (\x,0) {};
        \node[font=\small\bfseries\sffamily,text=Primary] at (\x,0.82) {\title};
      }
      \node[geom point] at (-5.2,-0.05) {};
      \node[geom label,above=2pt] at (-5.2,-0.05) {$A$};

      \draw[geom segment] (-1.55,-0.05)--(1.55,-0.05);
      \node[geom point] at (-1.55,-0.05) {};
      \node[geom label,above=2pt] at (-1.55,-0.05) {$A$};
      \draw[decorate,decoration={brace,mirror,amplitude=4pt},draw=Accent] (-1.55,-0.55)--(1.55,-0.55)
        node[midway,below=5pt,font=\scriptsize\sffamily,text=Accent] {$3\text{ см }1\text{ мм}$};

      \draw[geom segment] (3.65,-0.05)--(6.75,-0.05);
      \node[geom point] at (3.65,-0.05) {};
      \node[geom point] at (6.75,-0.05) {};
      \node[geom label,above=2pt] at (3.65,-0.05) {$A$};
      \node[geom label,above=2pt] at (6.75,-0.05) {$B$};
    \end{tikzpicture}
  \end{center}
}

\newcommand{\equalcontrastingsegmentsfigure}{%
  \begin{center}
    \begin{tikzpicture}[x=1cm,y=1cm]
      \coordinate (A) at (-6.6,1.0); \coordinate (B) at (-2.6,1.0);
      \coordinate (C) at (0.2,1.0);  \coordinate (D) at (4.2,1.0);
      \coordinate (M) at (-5.6,-0.7); \coordinate (N) at (1.0,-0.7);
      \draw[geom segment] (A)--(B); \draw[geom segment] (C)--(D); \draw[geom segment] (M)--(N);
      \foreach \P/\pos in {A/above,B/above,C/above,D/above,M/above,N/above}{\node[geom point] at (\P) {};\node[geom label,\pos=2pt] at (\P) {$\P$};}
      \node[geom caption,below=6pt] at (-4.6,1.0) {$AB=4\text{ см}$};
      \node[geom caption,below=6pt] at (2.2,1.0) {$CD=4\text{ см}$};
      \node[geom caption,below=6pt] at (-2.3,-0.7) {$MN=6\text{ см }6\text{ мм}$};
    \end{tikzpicture}
  \end{center}
}

\newcommand{\partssegmentfigure}{%
  \begin{center}
    \begin{tikzpicture}[x=1cm,y=1cm]
      \coordinate (A) at (-4.8,0); \coordinate (C) at (-0.8,0); \coordinate (B) at (4.2,0);
      \draw[geom segment] (A)--(B);
      \foreach \P/\pos in {A/above,C/above,B/above}{\node[geom point] at (\P) {};\node[geom label,\pos=3pt] at (\P) {$\P$};}
      \draw[decorate,decoration={brace,mirror,amplitude=4pt},draw=Accent] (A)+(0,-0.42)--($(C)+(0,-0.42)$)
        node[midway,below=5pt,font=\small\sffamily,text=Accent] {$AC$};
      \draw[decorate,decoration={brace,mirror,amplitude=4pt},draw=Secondary] (C)+(0,-0.42)--($(B)+(0,-0.42)$)
        node[midway,below=5pt,font=\small\sffamily,text=Secondary] {$CB$};
      \node[geom object label,above=10pt] at (-0.3,0) {$AB=AC+CB$};
    \end{tikzpicture}
  \end{center}
}

\newcommand{\partsproblemfigure}{%
  \begin{center}
    \begin{tikzpicture}[x=1cm,y=1cm]
      \coordinate (A) at (-4.5,0); \coordinate (C) at (-0.5,0); \coordinate (B) at (4.5,0);
      \draw[geom segment] (A)--(B);
      \foreach \P/\pos in {A/above,C/above,B/above}{\node[geom point] at (\P) {};\node[geom label,\pos=3pt] at (\P) {$\P$};}
      \draw[decorate,decoration={brace,mirror,amplitude=4pt},draw=Accent] (A)+(0,-0.45)--($(C)+(0,-0.45)$)
        node[midway,below=5pt,font=\small\sffamily,text=Accent] {$4\text{ см }8\text{ мм}$};
      \draw[decorate,decoration={brace,amplitude=4pt},draw=Secondary] (A)+(0,0.48)--($(B)+(0,0.48)$)
        node[midway,above=5pt,font=\small\sffamily,text=Secondary] {$9\text{ см }5\text{ мм}$};
      \node[geom caption,below=11pt] at (2.0,0) {$CB=?$};
    \end{tikzpicture}
  \end{center}
}

\newcommand{\studentconstructionfigure}{%
  \begin{center}
    \begin{tikzpicture}[x=1cm,y=1cm]
      \node[draw=GrayLine,line width=0.6pt,minimum width=14.7cm,minimum height=5.1cm] at (0,0) {};
      \coordinate (A) at (-5.4,1.25);
      \coordinate (C) at (-5.4,-1.35);
      \node[geom point] at (A) {}; \node[geom label,above=3pt] at (A) {$A$};
      \node[geom point] at (C) {}; \node[geom label,above=3pt] at (C) {$C$};
      \node[geom caption,anchor=west] at (-4.85,1.25) {постройте $AB=5\text{ см }8\text{ мм}$};
      \node[geom caption,anchor=west] at (-4.85,-1.35) {постройте $CD=3\text{ см }5\text{ мм}$};
    \end{tikzpicture}
  \end{center}
}

\newcommand{\studentconstructionanswersfigure}{%
  \begin{center}
    \begin{tikzpicture}[x=1cm,y=1cm]
      \node[draw=GrayLine,line width=0.6pt,minimum width=14.7cm,minimum height=5.1cm] at (0,0) {};
      \coordinate (A) at (-5.4,1.25); \coordinate (B) at (0.4,1.25);
      \coordinate (C) at (-5.4,-1.35); \coordinate (D) at (-1.9,-1.35);
      \draw[geom segment] (A)--(B); \draw[geom segment] (C)--(D);
      \foreach \P/\pos in {A/above,B/above,C/above,D/above}{\node[geom point] at (\P) {};\node[geom label,\pos=3pt] at (\P) {$\P$};}
      \node[geom caption,anchor=west] at (1.0,1.25) {$AB=5\text{ см }8\text{ мм}$};
      \node[geom caption,anchor=west] at (-1.3,-1.35) {$CD=3\text{ см }5\text{ мм}$};
    \end{tikzpicture}
  \end{center}
}

\newcommand{\readrulerpracticefigure}{%
  \begin{center}
    \begin{tikzpicture}[x=1cm,y=1cm]
      \node[geom card,minimum width=15.6cm,minimum height=5.15cm] at (0,0.55) {};
      \begin{scope}[shift={(-4,1.05)}]
        \drawruler[8]
        \coordinate (A) at (0,1.18); \coordinate (B) at (4.6,1.18);
        \draw[geom segment] (A)--(B);
        \node[geom point] at (A) {}; \node[geom point] at (B) {};
        \node[geom label,above=2pt] at (A) {$A$}; \node[geom label,above=2pt] at (B) {$B$};
        \draw[geom auxiliary,densely dashed] (A)--(0,0.45); \draw[geom auxiliary,densely dashed] (B)--(4.6,0.45);
      \end{scope}
      \begin{scope}[shift={(-4,-1.65)}]
        \drawruler[8]
        \coordinate (C) at (1.3,1.18); \coordinate (D) at (7.1,1.18);
        \draw[geom segment] (C)--(D);
        \node[geom point] at (C) {}; \node[geom point] at (D) {};
        \node[geom label,above=2pt] at (C) {$C$}; \node[geom label,above=2pt] at (D) {$D$};
        \draw[geom auxiliary,densely dashed] (C)--(1.3,0.45); \draw[geom auxiliary,densely dashed] (D)--(7.1,0.45);
      \end{scope}
    \end{tikzpicture}
  \end{center}
}
